\documentclass[10pt,a4paper]{amsart}
\usepackage[margin=19mm]{geometry}
\usepackage{amsmath,amssymb,mathtools}
\usepackage[scr=rsfs,cal=euler]{mathalfa}
\usepackage{xfrac}
\usepackage[unicode,hidelinks,bookmarks=false]{hyperref}
\newcommand{\ie}{i.e.\ }
\newcommand{\cf}{cf.\ }
\newcommand{\resp}{respectively\ }
\newcommand{\Subsection}{\subsection}
\providecommand{\nobreakdash}{\nobreak\hbox{-}}
\setlength{\emergencystretch}{2em}
\multlinegap0pt
\usepackage[shortlabels]{enumitem}
\usepackage{amssymb}
\usepackage{mathtools}
\newtheorem{lemma}{Lemma}
\newtheorem{proposition}{Proposition}
\title{Riccati Reductions for Modified Bessel Ratios:\\ Bernstein Positivity, Exact Certificates, and Transfer Obstructions}
\author{Domingos S. P. Salazar}
\date{Source: arXiv:2607.05538v1 (CC BY 4.0)}
\begin{document}
\maketitle

\noindent\emph{Source and licence.} Riccati Reductions for Modified Bessel Ratios:\\ Bernstein Positivity, Exact Certificates, and Transfer Obstructions. Version arXiv:2607.05538v1 (CC BY 4.0), distributed by arXiv under the Creative Commons Attribution 4.0 International licence, http://creativecommons.org/licenses/by/4.0/, as recorded in the arXiv licence metadata for this exact version and shown on the abstract page https://arxiv.org/abs/2607.05538v1. Full licence notice: \texttt{licenses/arxiv-2607.05538-CC-BY-4.0.txt}.\par\smallskip

\noindent\textit{Editorial scope.} Counted unit: the article's proposition prop:certificate-arithmetic (source0.tex lines 927--934). The article's complete original proof of it is delivered with it (source0.tex lines 935--979, one \texttt{proof} environment of 45 lines). No mathematics is written, repaired or re-derived: the statement and the proof are the article's own lines, in the article's own notation, and no retained line is substituted or re-flowed.  Retained uncounted as the source-local context the proof needs: lines 291--293; lines 912--920.  Nothing was omitted from any retained span, so the package carries no omission record.  No retained line contains a citation, so the wrapper carries no bibliography and no bibliography entry.  No retained line contains a figure environment, an image-inclusion command or drawing code, so no image is delivered and the package declares no asset dependency.  Editorial formatting only, in addition: an AMS \texttt{amsart} preamble that restates the article's own declarations and macros used by the retained lines, namely the environments \texttt{lemma}, \texttt{proposition} on separate counters, together with the standard packages those lines need.  The retained environments print in this document as Lemma 1, Proposition 1 in order of appearance, whereas the article prints the same items with the numbering of its own full document.  The complete native source of the article as distributed in the arXiv e-print for this exact version is bundled unchanged as \texttt{sources/204547/source0.tex}.
\begin{equation}\label{eq:I-series}
I_\nu(z)=\sum_{k\ge0}\frac{(z/2)^{2k+\nu}}{k!\,\Gamma(k+\nu+1)},
\end{equation}

\begin{lemma}[Geometric tail enclosure]\label{lem:tail}
Let \(t_k>0\) and suppose \(t_{k+1}/t_k\le q<1\) for all \(k\ge N+1\). Then
\begin{equation}\label{eq:tail-enclosure}
\sum_{k=0}^{N}t_k<\sum_{k=0}^{\infty}t_k
\le\sum_{k=0}^{N}t_k+\frac{t_{N+1}}{1-q}.
\end{equation}
If the \(t_k\) and \(q\) are rational, both enclosing bounds are explicit
rationals.
\end{lemma}

\begin{proposition}[Exact arithmetic package]\label{prop:certificate-arithmetic}
The following two rational statements hold:
\begin{enumerate}[label=\textup{(\roman*)},leftmargin=2.3em]
\item \(I_1(10)/I_0(10)<9487/10000\).
\item If \(h(r)=199/200-r/10-r^2\), then
\(h(9487/10000)=9831/10^8\).
\end{enumerate}
\end{proposition}

\begin{proof}
From \eqref{eq:I-series},
\begin{equation}\label{eq:appendix-bessel-series}
I_0(10)=\sum_{k=0}^{\infty}\frac{25^k}{(k!)^2},
\qquad
I_1(10)=\sum_{k=0}^{\infty}\frac{5\cdot25^k}{k!\,(k+1)!}.
\end{equation}
Both series have positive rational terms. Set
\begin{equation}\label{eq:appendix-L0}
L_0=\sum_{k=0}^{20}\frac{25^k}{(k!)^2}<I_0(10).
\end{equation}
For the \(I_1\) terms, for \(k\ge21\),
\begin{equation}\label{eq:appendix-tail-ratio}
\frac{t_{k+1}}{t_k}
=\frac{25}{(k+1)(k+2)}
\le \frac{25}{22\cdot23}
=\frac{25}{506}.
\end{equation}
Therefore Lemma~\ref{lem:tail} gives
\begin{equation}\label{eq:appendix-U1}
I_1(10)\le U_1:=
\sum_{k=0}^{20}\frac{5\cdot25^k}{k!\,(k+1)!}
+\frac{5\cdot25^{21}}{21!\,22!}\cdot\frac{506}{481}.
\end{equation}
The exact rational comparison is
\begin{equation}\label{eq:appendix-integer-witness}
9487L_0-10000U_1
=
\frac{9066027287067976494607015156102735792727}
{3214199764297248897429588860517482496}>0.
\end{equation}
Hence
\begin{equation}\label{eq:appendix-ratio-certificate}
\frac{I_1(10)}{I_0(10)}
<\frac{U_1}{L_0}
<\frac{9487}{10000},
\end{equation}
which proves (i). For (ii), direct rational evaluation gives
\begin{equation}\label{eq:appendix-h-value}
h\!\left(\frac{9487}{10000}\right)
=\frac{199}{200}-\frac{9487}{100000}
-\left(\frac{9487}{10000}\right)^2
=\frac{9831}{10^8}.
\end{equation}
\end{proof}

\end{document}
